% Part: first-order-logic
% Chapter: natural-deduction
% Section: provability-propositional

% pamriksan sipat-sipat provabilitas kang dibutuhake kanggo himpunan
% konsisten maksimal ing bab kelengkapan.

\documentclass[../../../include/open-logic-section]{subfiles}

\begin{document}

\iftag{FOL}
      {\olfileid{fol}{ntd}{ppr}}
      {\olfileid{pl}{ntd}{ppr}}

\olsection{\usetoken{S}{derivability} lan Konektif Proposisional}

\begin{explain}
  Kita netepake yen relasi !!{derivability}~$\Proves$ saka deduksi
  natural cukup kuwat kanggo netepake sawetara kasunyatan dhasar kang
  ngemot konektif proposisional, kayata $!A \land !B \Proves !A$ lan
  $!A, !A \lif !B \Proves !B$ (modus ponens). Kasunyatan kasebut
  dibutuhake ing bukti teorema kelengkapan.
\end{explain}

\begin{prop}\ollabel{prop:provability-land}
  \begin{enumerate}
  \item \ollabel{prop:provability-land-left} $!A \land !B \Proves !A$
    lan $!A \land !B \Proves !B$ kalorone lumaku.
  \item \ollabel{prop:provability-land-right} $!A, !B \Proves !A \land !B$.
  \end{enumerate}
\end{prop}

\begin{proof}
  \begin{enumerate}
  \item Kita bisa !!{derive} kalorone:
    \begin{prooftree}
      \AxiomC{$!A \land !B$}
      \RightLabel{\Elim{\land}}
      \UnaryInfC{$!A$}
      \DisplayProof\qquad\bottomAlignProof
      \AxiomC{$!A \land !B$}
      \RightLabel{\Elim{\land}}
      \UnaryInfC{$!B$}
    \end{prooftree}
  \item Kita bisa !!{derive}:
    \begin{prooftree}
      \AxiomC{$!A$}
      \AxiomC{$!B$}
      \RightLabel{\Intro{\land}}
      \BinaryInfC{$!A \land !B$}
    \end{prooftree}
  \end{enumerate}
\end{proof}


\begin{prop}\ollabel{prop:provability-lor}
  \begin{enumerate}
  \item $!A \lor !B, \lnot !A, \lnot !B$ ora konsisten.
  \item $!A \Proves !A \lor !B$ lan $!B \Proves !A \lor !B$
    kalorone lumaku.
  \end{enumerate}
\end{prop}

\begin{proof}
  \begin{enumerate}
  \item Gatekna !!{derivation} iki:
    \begin{prooftree}
      \AxiomC{$!A \lor !B$}
      \AxiomC{$\lnot !A$}
      \AxiomC{$\Discharge{!A}{1}$}
      \RightLabel{\Elim{\lnot}}
      \BinaryInfC{$\lfalse$}
      \AxiomC{$\lnot !B$}
      \AxiomC{$\Discharge{!B}{1}$}
      \RightLabel{\Elim{\lnot}}
      \BinaryInfC{$\lfalse$}
      \DischargeRule{\Elim{\lor}}{1}
      \TrinaryInfC{$\lfalse$}
    \end{prooftree}
    Iki !!a{derivation} saka~$\lfalse$ adhedhasar asumsi
    !!{undischarged} $!A \lor !B$, $\lnot !A$, lan $\lnot !B$.
  \item Kita bisa !!{derive} kalorone:
    \begin{prooftree}
      \AxiomC{$!A$}
      \RightLabel{\Intro{\lor}}
      \UnaryInfC{$!A \lor !B$}
      \DisplayProof\qquad\bottomAlignProof
      \AxiomC{$!B$}
      \RightLabel{\Intro{\lor}}
      \UnaryInfC{$!A \lor !B$}
    \end{prooftree}
  \end{enumerate}
\end{proof}

\begin{prop}\ollabel{prop:provability-lif}
  \begin{enumerate}
  \item \ollabel{prop:provability-lif-left} $!A, !A \lif !B \Proves !B$.
  \item \ollabel{prop:provability-lif-right}
    $\lnot !A \Proves !A \lif !B$ lan $!B \Proves !A \lif !B$
    kalorone lumaku.
  \end{enumerate}
\end{prop}

\begin{proof}
  \begin{enumerate}
  \item Kita bisa !!{derive}:
    \begin{prooftree}
      \AxiomC{$!A \lif !B$}
      \AxiomC{$!A$}
      \RightLabel{\Elim{\lif}}
      \BinaryInfC{$!B$}
    \end{prooftree}

  \item Iki dituduhake dening loro !!{derivation}s iki:
    \begin{prooftree}
      \AxiomC{$\lnot !A$}
      \AxiomC{$\Discharge{!A}{1}$}
      \RightLabel{\Elim{\lnot}}
      \BinaryInfC{$\lfalse$}
      \RightLabel{\FalseInt}
      \UnaryInfC{$!B$}
      \DischargeRule{\Intro{\lif}}{1}
      \UnaryInfC{$!A \lif !B$}
      \DisplayProof\qquad\bottomAlignProof
      \AxiomC{$!B$}
      \RightLabel{\Intro{\lif}}
      \UnaryInfC{$!A \lif !B$}
    \end{prooftree}
    Gatekna yen $\Intro{\lif}$ bisa, nanging ora kudu,
    !!{discharge} asumsi~$!A$.
  \end{enumerate}
\end{proof}

\end{document}
